\section{Matrix-free Z-vectors and bounded-memory response}
\label{sec:matrix-free-bounded-memory}

\subsection{Response dimension and dense-matrix cost}

At k-point $\vk$, occupied--virtual rotations have complex dimension
$n_o^{\vk}n_v^{\vk}$.  After the realification of
Eq.~\eqref{eq:realified-coordinate}, the total real dimension is
\begin{equation}
  m=2\sum_{\vk}n_o^{\vk}n_v^{\vk}.
  \label{eq:response-real-dimension}
\end{equation}
A dense real Hessian is $m\times m$ and its float64 storage alone is
\begin{equation}
  M_{\mathrm{dense}}=8m^2\ \text{bytes}.
  \label{eq:dense-hessian-storage}
\end{equation}
Factorization and condition diagnostics require additional copies and
workspace.  The dense route is useful for equation-level parity on small
reference systems, but it is not a production scaling strategy.

Each dense column requires one AX application to a basis vector $e_j$.  This
does not mean that $m^2$ separate equations are solved; rather, $m$ operator
actions are collected to materialize $m^2$ matrix elements.

\subsection{Matrix-free Hessian actions}

The matrix-free route never stores $A$.  For a trial vector $u$, it evaluates
\begin{equation}
  u\longmapsto D^{(1)}(u)
  \longmapsto \delta J[D^{(1)}]-\tfrac12\delta K[D^{(1)}]
  \longmapsto Au.
  \label{eq:matrix-free-action}
\end{equation}
The first map constructs a Hermitian transition density from VO rotations.
The second applies the selected FFTDF/GDF/RSDF central J/K provider.  The last
projects the AO result back to VO blocks and adds the orbital-energy gap.

If a Krylov solve takes 21 iterations, it does not select only 21 important
components out of the $m$ unknowns.  Every iteration applies AX to a complete
length-$m$ trial vector; 21 Krylov directions were sufficient to represent
the solution to the requested residual.

\subsection{MINRES and the energy-gap preconditioner}

The realified HF Hessian is ideally real symmetric, but need not be positive
definite.  MINRES is a Lanczos/Krylov method that reduces the residual norm
for symmetric indefinite systems.  Starting from $z_0=0$, the order-$p$
subspace is
\begin{equation}
  \mathcal K_p(A,b)
  =\operatorname{span}\{b,Ab,A^2b,\ldots,A^{p-1}b\}.
  \label{eq:minres-krylov}
\end{equation}

A simple preconditioner is the independent-particle gap,
\begin{equation}
  P_{ai\vk,ai\vk}=\varepsilon_a^{\vk}-\varepsilon_i^{\vk},
  \qquad P^{-1}r\approx A^{-1}r.
  \label{eq:gap-preconditioner}
\end{equation}
This does not approximate the final physical solution by dropping
two-electron coupling.  MINRES still applies the full $A$; $P^{-1}$ rescales
poorly conditioned directions so the Krylov subspace converges faster.
Positive VO gaps, dtype, and k-point ordering are checked before iteration.

Convergence is accepted using the explicitly recomputed equation residual
\begin{equation}
  \eta_z
  =\frac{\|A_R^Tz_R+g_R\|_\infty}
         {\max(1,\|g_R\|_\infty)}.
  \label{eq:scaled-z-residual}
\end{equation}
If needed, iterative refinement solves
\begin{equation}
  A_R^T\delta z_R=-\left(A_R^Tz_R+g_R\right),
  \qquad z_R\leftarrow z_R+\delta z_R
  \label{eq:iterative-refinement}
\end{equation}
with the same matrix-free operator.

\subsection{Symmetry, complex k-points, and solver selection}

MINRES assumes that the realified operator is symmetric within numerical
tolerance.  Incorrect reverse-pair conjugation, $1/N_k$, occupation factors,
or J/K transposes increase $\|A_R-A_R^T\|$ in a dense test and can destabilize
the matrix-free solve.  Small Gamma and sampled-k systems must establish
dense/matrix-free parity for solutions, gradients, and residuals.

An automatic policy may use dense reference solves at small $m$ and
matrix-free solves at large $m$.  Rebuilding a large dense matrix merely to
obtain an exact condition number defeats the purpose.  Production reports
should state that the condition number was not evaluated and preserve the
iteration count, residual, and minimum gap instead.

\subsection{Batching nuclear sources}

The Z-vector is solved once, but $B^{Ax}$ contains $3N_A$ perturbations.
Materializing all AO/VO response tensors together scales approximately as
\begin{equation}
  M_B\sim 3N_A N_kN_{\mathrm{AO}}^2\,s_{\mathrm{dtype}}.
  \label{eq:nuclear-batch}
\end{equation}
For a component batch $\mathcal C$, one evaluates
\begin{equation}
  g_{Ax}^{\mathrm{resp}}=z^TB^{Ax},\qquad (A,x)\in\mathcal C,
\end{equation}
and keeps only the scalar gradient.  The batch size is determined from free
memory, dtype, AO dimension, and temporary count.  An OOM retry starts from
the same component with a smaller batch and must not add completed
components twice.

\subsection{Adjoint contractions instead of full derivative tensors}

Let a whitened three-center factor be $C=L^{-1}B$.  Its nuclear derivative is
\begin{equation}
  C^{Ax}=L^{-1}\left(B^{Ax}-L^{Ax}C\right),
  \label{eq:bounded-cderi-derivative}
\end{equation}
but the gradient usually needs only
\begin{equation}
  \mathcal G^{Ax}=\operatorname{Re}\langle W_C,C^{Ax}\rangle.
  \label{eq:reverse-contraction}
\end{equation}
Defining $W_B=L^{-\dagger}W_C$ gives
\begin{equation}
  \mathcal G^{Ax}
  =\operatorname{Re}\langle W_B,B^{Ax}\rangle
   -\operatorname{Re}\langle W_B,L^{Ax}C\rangle.
  \label{eq:reverse-cderi-response}
\end{equation}
Source and metric derivative blocks can therefore be accumulated directly
into scalars without storing a full
$(3N_A,N_{\mathrm{aux}},N_{\mathrm{AO}}^2)$ tensor.  This is reverse-mode or
adjoint algebra, not a screening approximation.

\subsection{Signed two-density polarization}

Write the fixed-density restricted-HF two-electron derivative as
\begin{equation}
  Q^{Ax}(X)
  =\frac{1}{2N_k}\sum_{\vk}\operatorname{Re}
   \Tr\left[X^{\vk}L^{Ax}[X]^{\vk}\right].
  \label{eq:quadratic-derivative}
\end{equation}
Because $L^{Ax}$ is linear in density, the response between reference density
$D$ and Z transition density $T$ obeys the exact polarization identity
\begin{equation}
  R_{2e}^{Z,Ax}=Q^{Ax}(D+T)-Q^{Ax}(D-T).
  \label{eq:polarization-identity}
\end{equation}
The equivalent direct form is
\begin{equation}
  R_{2e}^{Z,Ax}
  =\frac1{N_k}\sum_{\vk}\operatorname{Re}\left\{
    \Tr[D^{\vk}L^{Ax}[T]^{\vk}]
   +\Tr[T^{\vk}L^{Ax}[D]^{\vk}]
  \right\}.
  \label{eq:bilinear-response}
\end{equation}

Processing the signed states in one SR/LR/metric traversal can reuse grids
and integrals that two independent quadratic traversals would read twice.
Both Coulomb and exchange orientations, auxiliary-metric response, the Ewald
exchange correction, and k-pair conjugation must be retained.  Evaluating
only $D+T$ mixes self terms and is not Eq.~\eqref{eq:bilinear-response}.

\subsection{Streaming, block size, and object lifetime}

If a reciprocal derivative treats $b_G$ vectors and all AO pairs together,
the main working set scales as
\begin{equation}
  M_{\mathrm{block}}
  \sim b_GN_kN_{AO}^2s_{\mathrm{complex}}\,n_{\mathrm{temporary}}.
  \label{eq:block-working-set}
\end{equation}
Reducing $b_G$ lowers peak memory while retaining every reciprocal vector.
Smaller blocks may add kernel launches, so a free-memory fraction and a hard
cap are used together and the resolved block size and count are reported.

CDERI, multigrid tasks, and response temporaries should also be released
after their last consumer.  The central distinction is
\begin{equation}
  \underbrace{\text{batching, streaming, adjoint, lifetime}}_{\text{same finite representation}}
  \neq
  \underbrace{\text{looser screening, smaller mesh, smaller basis}}_{\text{changed representation}}.
  \label{eq:algebra-vs-approx}
\end{equation}
The first group must agree within roundoff and is tested by parity.  The
second changes both cost and accuracy and requires a convergence study.

\subsection{Algorithmic references}

MINRES was introduced by C.~C. Paige and M.~A. Saunders,
\href{https://doi.org/10.1137/0712047}{SIAM J. Numer. Anal. 12, 617--629
(1975)}.  This project combines that general algorithm with an HF Hessian
action and orbital-gap preconditioner.  The PySCF design and periodic Gaussian
framework are described by
\href{https://doi.org/10.1002/wcms.1340}{Sun et al., WIREs Comput. Mol. Sci.
8, e1340 (2018)}.
